\begin{lemma}[Identische Lyman-Serienverschiebung und differentieller Lamb-Shift-Offset]
\label{lem:lyman_universal_shift}
Sei $\widehat{H}_{\rm int} = \eta_A^{\rm eff} \, \widehat{V}$ der effektive isotrope Nahfeld-Störoperator erster Ordnung im EABC-Modell mit $\eta_A^{\rm eff} = \frac{4}{7}\eta_A$ und dem unpolarisierten Kontaktterm $\widehat{V} = \delta^3(\mathbf{r})$ in atomaren Standardeinheiten ($E_{\rm Hartree} = 2 R_\infty h c$, $a_0 = 1$). 
Für den oberen Schrankenwert der Fisher-Informationsmatrix $\eta_A = 1.06 \times 10^{-14}$ ($\eta_A^{\rm eff} \approx 6.057 \times 10^{-15}$) gilt:

\begin{enumerate}[(i)]
    \item \textbf{Invarianz der Lyman-Emissionsfrequenzen:} Alle atomaren Übergänge der Lyman-Serie $\text{Ly}_n$ ($nP_{1/2} \to 1S_{1/2}$ für $n \ge 2$) erfahren eine identische negative Frequenzverschiebung:
    \begin{equation}
        \Delta\nu(\text{Ly}_n) \equiv \Delta\nu_{\rm Ly} = -\frac{\eta_A^{\rm eff}}{h} \langle 1S | \delta^3(\mathbf{r}) | 1S \rangle = -39.85\,\text{Hz} \quad \forall \, n \ge 2.
    \end{equation}
    
    \item \textbf{Differentieller $1S\text{--}2S$-Offset und Lamb-Shift-Identität:} Die Frequenzdifferenz zwischen dem Doppler-freien Zwei-Photonen-Kanal ($1S_{1/2} \to 2S_{1/2}$) und der fundamentalen Lyman-$\alpha$-Resonanz ($2P_{1/2} \to 1S_{1/2}$) ist strikt positiv und identisch mit der Netto-Verschiebung der klassischen Mikrowellen-Lamb-Verschiebung $\Delta\nu_{\rm Lamb}(2S_{1/2} - 2P_{1/2})$:
    \begin{equation}
        \delta\Delta\nu \coloneqq \Delta\nu(1S \to 2S) - \Delta\nu(\text{Ly}_\alpha) = \Delta\nu_{\rm Lamb}(2S_{1/2} - 2P_{1/2}) = +4.98\,\text{Hz}.
    \end{equation}
\end{enumerate}
\end{lemma}

\begin{proof}
Zu (i): Die Frequenzverschiebung eines optischen Übergangs $|i\rangle \to |f\rangle$ ergibt sich in erster Ordnung Störungstheorie zu:
\begin{equation}
    \Delta\nu_{i \to f} = \frac{\Delta E_i - \Delta E_f}{h} = \frac{\eta_A^{\rm eff}}{h} \left( \langle i | \delta^3(\mathbf{r}) | i \rangle - \langle f | \delta^3(\mathbf{r}) | f \rangle \right).
\end{equation}
Für nicht-relativistische Coulomb-Eigenfunktionen gilt am Koordinatenursprung $\mathbf{r} = \mathbf{0}$:
\begin{equation}
    |\psi_{n\ell m}(\mathbf{0})|^2 = \frac{1}{\pi a_0^3 n^3} \, \delta_{\ell, 0}.
\end{equation}
Folglich verschwindet der Erwartungswert des Kontaktterms für alle Bahndrehimpulszustände mit $\ell \ge 1$ identisch:
\begin{equation}
    \langle nP_{1/2} | \delta^3(\mathbf{r}) | nP_{1/2} \rangle = 0 \quad \forall \, n \ge 2.
\end{equation}
Für den Grundzustand $|1S\rangle$ ($n=1, \ell=0$) lautet das Matrixelement hingegen:
\begin{equation}
    \langle 1S | \delta^3(\mathbf{r}) | 1S \rangle = \frac{1}{\pi a_0^3}.
\end{equation}
Daraus folgt unmittelbar für alle Linien der Lyman-Serie ($nP_{1/2} \to 1S_{1/2}$):
\begin{equation}
    \Delta\nu(\text{Ly}_n) = \frac{\eta_A^{\rm eff}}{h} \left( 0 - \frac{1}{\pi a_0^3} \right) = -\eta_A^{\rm eff} \, \frac{2 R_\infty c}{1^3} = -\eta_A^{\rm eff} \, (2 R_\infty c).
\end{equation}
Mit $2 R_\infty c \approx 6.579\,684 \times 10^{15}\,\text{Hz}$ und $\eta_A^{\rm eff} = \frac{4}{7}(1.06 \times 10^{-14}) \approx 6.057 \times 10^{-15}$ ergibt sich:
\begin{equation}
    \Delta\nu(\text{Ly}_n) = - (6.057 \times 10^{-15}) \times (6.579\,684 \times 10^{15}\,\text{Hz}) \approx -39.853\,\text{Hz}.
\end{equation}

Zu (ii): Für den metastabilen Zustand $|2S\rangle$ ($n=2, \ell=0$) beträgt die Aufenthaltswahrscheinlichkeit am Ursprung:
\begin{equation}
    \langle 2S | \delta^3(\mathbf{r}) | 2S \rangle = \frac{1}{\pi a_0^3 \cdot 2^3} = \frac{1}{8\pi a_0^3}.
\end{equation}
Die Frequenzverschiebung des Zwei-Photonen-Übergangs $1S \to 2S$ errechnet sich folglich zu:
\begin{equation}
    \Delta\nu(1S \to 2S) = \frac{\eta_A^{\rm eff}}{h} \left( \frac{1}{8\pi a_0^3} - \frac{1}{\pi a_0^3} \right) = -\frac{7}{8} \eta_A^{\rm eff} (2 R_\infty c) = \frac{7}{8} \Delta\nu_{\rm Ly} \approx -34.871\,\text{Hz}.
\end{equation}
Subtraktion von $\Delta\nu(\text{Ly}_\alpha)$ liefert:
\begin{equation}
    \delta\Delta\nu = \Delta\nu(1S \to 2S) - \Delta\nu(\text{Ly}_\alpha) = -\frac{7}{8} \Delta\nu_{\rm Ly} - \Delta\nu_{\rm Ly} = +\frac{1}{8} |\Delta\nu_{\rm Ly}| = +4.982\,\text{Hz}.
\end{equation}
Für die Mikrowellen-Lamb-Verschiebung $\nu_{\rm Lamb} = \nu(2S_{1/2} - 2P_{1/2})$ gilt andererseits:
\begin{equation}
    \Delta\nu_{\rm Lamb} = \frac{\eta_A^{\rm eff}}{h} \left( \langle 2S | \delta^3(\mathbf{r}) | 2S \rangle - \langle 2P | \delta^3(\mathbf{r}) | 2P \rangle \right) = \frac{\eta_A^{\rm eff}}{h} \left( \frac{1}{8\pi a_0^3} - 0 \right) = +\frac{1}{8} |\Delta\nu_{\rm Ly}| \approx +4.982\,\text{Hz}.
\end{equation}
Damit ist $\delta\Delta\nu \equiv \Delta\nu_{\rm Lamb} = +4.98\,\text{Hz}$ exakt bewiesen.
\end{proof}

\begin{remark}[Spektroskopische Falsifikations-Signatur]
Das Lemma liefert das exakte algebraische Fundament für die Brechung der Kollinearität in Laborpfad~I: Während die globale Planck-Skalen-Variation $\epsilon_h$ über die Rydberg-Konstante $\nu \propto h^{-3}$ alle optischen Übergänge ($1S\text{--}2S$ und $\text{Ly}_\alpha$) mit nahezu identischer Steigung skaliert ($\partial_{\epsilon_h}\nu \approx -3\nu$), ist die differentielle Differenz $\delta\Delta\nu$ immun gegen reine Coulomb-Gradienten und isoliert den oktonionischen Nahfeld-Term $\eta_A^{\rm eff}$ im Frequenzraum mit einer Signatur von $+4.98\,\text{Hz}$.
\end{remark}